% 2007 年广东卷理科第 6 题：图 2 算法流程图。
%   开始 → 输入 A1,A2,…,A10 → s=0,i=4 → 判断框（题目待填条件）→
%   是 → s=s+A_i → i=i+1 → 回到判断框上方；
%   否 → 输出 s → 结束。
% 来源 PDF：gaokaomath-src/普通高考/2007/2007广东理.pdf
% SHA-256：655db04340f14b909d7f248531436d45dd8cbd71e69aca976c3bd2d1e568ceed
% 原页：2007广东理.pdf 第 1 页（A3 三栏）第 2 栏，书页码 371。
% 重绘依据：按原页流程与分支方向重绘；判断框内留空，对应题目要求填写的条件。
\begin{tikzpicture}[scale=1.0, semithick,
  every node/.style={inner sep=2pt},
  box/.style={draw, rounded corners=2pt, minimum width=2.1cm,
    minimum height=0.55cm, align=center},
  para/.style={draw, trapezium, trapezium left angle=70,
    trapezium right angle=110, minimum width=2.9cm, minimum height=0.55cm,
    align=center},
  dec/.style={draw, diamond, aspect=2.6, minimum width=2.8cm,
    minimum height=1.05cm, align=center},
  arr/.style={->, >=stealth}]

  \node[box]  (start) at (0,8.6) {开始};
  \node[para] (in)    at (0,7.8) {输入 $A_1,A_2,\cdots,A_{10}$};
  \node[box]  (init)  at (0,6.9) {$s=0,\ i=4$};
  \node[dec]  (test)  at (0,5.5) {};
  \node[box]  (add)   at (4.6,4.9) {$s=s+A_i$};
  \node[box]  (inc)   at (4.6,6.1) {$i=i+1$};
  \node[para] (out)   at (0,3.9) {输出 $s$};
  \node[box]  (end)   at (0,3.0) {结束};

  \draw[arr] (start) -- (in);
  \draw[arr] (in) -- (init);
  \draw[arr] (init) -- (test);
  \draw[arr] (test.east) -- node[above] {是} (add.west);
  \draw[arr] (add.north) -- (inc.south);
  \draw[arr] (inc.west) -- (0,6.1) -- (0,6.32);
  \draw[arr] (test.south) -- node[left] {否} (out.north);
  \draw[arr] (out) -- (end);

  \node[below] at (2.6,2.55) {图 2};
\end{tikzpicture}
